\documentclass[aspectratio=169,10pt,brazil]{beamer}
\usepackage[brazil]{babel}

\usetheme{UELaula}

\title[Busca binária]{Busca binária}
\subtitle{Procurando em listas ordenadas}
\author[Prof. Nome do Professor]{Prof. Dr. Nome do Professor}
\institute{Departamento de Computação\\Universidade Estadual de Londrina}
\disciplina{Algoritmos e Estruturas de Dados}
\aula{5}
\semestre{2026/1}

\begin{document}

\maketitle

\begin{frame}{Roteiro}
  \tableofcontents
\end{frame}

\section{Ideia}

\begin{frame}{Dividir para conquistar}{Metade do trabalho a cada passo}
  \begin{itemize}
    \item Compare a chave com o elemento do \alert{meio}
    \item Descarte a metade que não pode conter a chave
    \item Repita na metade restante
  \end{itemize}
\end{frame}

\section{Análise}

\begin{frame}{Definição e teorema}
  \begin{definition}
    Uma lista $A[1..n]$ é \emph{ordenada} se $A[i] \le A[i+1]$ para todo $i < n$.
  \end{definition}
  \begin{theorem}
    A busca binária faz no máximo $\lfloor \log_2 n \rfloor + 1$ comparações.
  \end{theorem}
  \begin{proof}
    A cada passo o intervalo de busca cai pela metade, e $n$ só pode ser dividido
    por 2 $\lfloor \log_2 n \rfloor$ vezes antes de chegar a 1.
  \end{proof}
\end{frame}

\begin{frame}[fragile]{Implementação}
\begin{lstlisting}[language=Python]
def busca_binaria(a, chave):
    lo, hi = 0, len(a) - 1
    while lo <= hi:
        meio = (lo + hi) // 2      # ponto médio
        if a[meio] == chave:
            return meio
        if a[meio] < chave:
            lo = meio + 1
        else:
            hi = meio - 1
    return -1
\end{lstlisting}
\end{frame}

\begin{frame}{Exercícios}
  \begin{exercicio}
    Execute a busca binária pela chave 23 em $[2, 5, 8, 12, 16, 23, 38]$ e anote
    cada comparação.
  \end{exercicio}
  \begin{exercicio}[difícil]
    Mostre que o ponto médio $\lfloor (lo + hi)/2 \rfloor$ nunca sai de $[lo, hi]$.
  \end{exercicio}
\end{frame}

\begin{frame}{Comparação}
  \centering
  \begin{tabular}{lrr}
    \hline
    Método & Pior caso & Pré-condição \\ \hline
    Linear & $n$ & nenhuma \\
    Binária & $\log_2 n$ & lista ordenada \\ \hline
  \end{tabular}
\end{frame}

\end{document}
